\chapter{Phân tích yêu cầu hệ thống}
\label{ch:yeucau}

\section{Tác nhân}
\label{sec:tacnhan}

Hệ thống có bốn tác nhân:
\begin{enumerate}
    \item \textbf{Khách hàng (customer):} cá nhân hoặc đại diện doanh nghiệp thuê chỗ làm việc, phòng họp và tham gia cộng đồng.
    \item \textbf{Nhân viên quầy (staff):} làm việc tại quầy lễ tân của một chi nhánh; check-in, check-out, tạo đơn cho khách vãng lai, thu tiền mặt, ghi nhận dịch vụ phát sinh và khóa chỗ khi cần bảo trì.
    \item \textbf{Quản lý chi nhánh (branch admin):} phụ trách một chi nhánh; cấu hình mặt bằng, bảng giá, chính sách hủy, dịch vụ; quản lý nhân viên, xử lý hoàn tiền và xem báo cáo của chi nhánh.
    \item \textbf{Quản trị viên hệ thống (super admin):} quản lý toàn chuỗi; tạo chi nhánh, quản lý danh mục dùng chung, phân quyền, khuyến mãi, hạng thành viên và xem nhật ký kiểm toán.
\end{enumerate}

Khách hàng và quản trị viên hệ thống không gắn với chi nhánh (\mt{branch\_id} rỗng); nhân viên và quản lý chi nhánh bắt buộc gắn với đúng một chi nhánh. Quy tắc này được kiểm tra ở tầng ứng dụng và bằng ràng buộc \mt{check\_branch\_by\_role} trong cơ sở dữ liệu (Bảng~\ref{tab:tacnhan}).

\begin{table}[!htbp]
\centering
\caption{Tác nhân và phạm vi dữ liệu}
\label{tab:tacnhan}
\small
\renewcommand{\arraystretch}{1.2}
\begin{tabularx}{\linewidth}{|P{2.9cm}|P{2.2cm}|L|L|}
\hline
\textbf{Tác nhân} & \textbf{Chi nhánh} & \textbf{Phạm vi dữ liệu} & \textbf{Công việc chính} \\
\hline
Khách hàng (\mt{customer}) & Không gắn & Không gian đang mở; hồ sơ, đơn và giao dịch của chính mình & Tìm chỗ, đặt chỗ, thanh toán, hủy đơn, kết nối thành viên \\
\hline
Nhân viên (\mt{staff}) & Bắt buộc gắn & Đơn, check-in, dịch vụ phát sinh, bảo trì của chi nhánh mình & Vận hành quầy \\
\hline
Quản lý chi nhánh (\mt{branch\_admin}) & Bắt buộc gắn & Không gian, bảng giá, chính sách, nhân viên, hoàn tiền, báo cáo của chi nhánh mình & Cấu hình và giám sát chi nhánh \\
\hline
Quản trị hệ thống (\mt{super\_admin}) & Không gắn & Toàn bộ dữ liệu, danh mục dùng chung, nhật ký kiểm toán & Quản lý chuỗi, phân quyền \\
\hline
\end{tabularx}
\end{table}

\section{Yêu cầu chức năng}
\label{sec:yeucau_chucnang}

Các yêu cầu chức năng được đánh mã \mt{FR-xx}.

\subsection{Khách hàng}
\begin{itemize}
    \item \textbf{FR-01 (Đăng ký, đăng nhập):} đăng ký bằng email và mật khẩu hoặc đăng nhập bằng tài khoản Google qua Supabase Auth.
    \item \textbf{FR-02 (Hồ sơ năng lực):} cập nhật thông tin cá nhân, ảnh đại diện, nghề nghiệp, công ty, kỹ năng (kèm mức độ 1--5), sở thích và thông tin liên hệ; chọn công khai hoặc ẩn thông tin liên hệ.
    \item \textbf{FR-03 (Xem chi nhánh và sơ đồ):} xem danh sách chi nhánh, bảng giá, ảnh không gian và sơ đồ mặt bằng từng tầng.
    \item \textbf{FR-04 (Tìm chỗ):} lọc theo chi nhánh, tầng, loại không gian và khung giờ; xem tình trạng từng chỗ trên sơ đồ.
    \item \textbf{FR-05 (Đặt chỗ):} đặt theo giờ, ngày, tuần hoặc tháng; đơn mới ở trạng thái \mt{PENDING\_PAYMENT} và được giữ 15 phút; mỗi khách có tối đa 3 đơn chờ thanh toán cùng lúc.
    \item \textbf{FR-06 (Thanh toán trực tuyến):} thanh toán bằng mã VietQR qua PayOS hoặc ví MoMo; đơn chuyển sang \mt{CONFIRMED} khi hệ thống nhận được xác nhận thanh toán.
    \item \textbf{FR-07 (Dịch vụ thêm và ưu đãi):} chọn dịch vụ đi kèm khi đặt chỗ; áp dụng mã khuyến mãi hoặc voucher cá nhân; tự động được giảm giá theo hạng thành viên.
    \item \textbf{FR-08 (Lịch sử và mã đặt chỗ):} xem đơn sắp tới, đã hoàn thành và đã hủy; xem mã đặt chỗ và mã QR để check-in; gia hạn hoặc gọi thêm dịch vụ cho đơn đang hiệu lực.
    \item \textbf{FR-09 (Hủy đơn):} hủy đơn trước giờ bắt đầu; hệ thống tính số tiền hoàn và phí phạt theo chính sách hủy.
    \item \textbf{FR-10 (Gợi ý và kết nối đối tác):} xem danh sách thành viên phù hợp kèm điểm tương đồng và lý do gợi ý; gửi, chấp nhận hoặc từ chối lời mời kết nối.
    \item \textbf{FR-11 (Bảng tin cộng đồng):} đăng bài, gắn thẻ chủ đề và lọc bài theo thẻ.
    \item \textbf{FR-12 (Trợ lý ảo):} hỏi đáp bằng ngôn ngữ tự nhiên về chi nhánh, chỗ trống, giá, chính sách và đơn của mình; nhận đề xuất đặt chỗ, hủy chỗ hoặc gọi thêm dịch vụ và xác nhận trước khi thực hiện.
\end{itemize}

\subsection{Nhân viên quầy}
\begin{itemize}
    \item \textbf{FR-13 (Bảng điều khiển ca trực):} xem các chỉ số trong ngày của chi nhánh: số đơn, số khách đang có mặt, số chỗ trống, danh sách đơn sắp đến.
    \item \textbf{FR-14 (Đặt chỗ tại quầy):} tạo đơn cho khách vãng lai, gắn với tài khoản có sẵn hoặc tạo nhanh tài khoản khách.
    \item \textbf{FR-15 (Thu tiền mặt):} xác nhận đã thu tiền mặt; đơn chuyển sang \mt{CONFIRMED}.
    \item \textbf{FR-16 (Check-in):} nhập mã đặt chỗ hoặc quét QR; chỉ cho phép từ 30 phút trước giờ bắt đầu đến giờ kết thúc của đơn.
    \item \textbf{FR-17 (Check-out):} trả chỗ cho khách; hệ thống chặn check-out khi còn dịch vụ phát sinh chưa thanh toán hoặc phí trả chỗ muộn chưa được xử lý.
    \item \textbf{FR-18 (Dịch vụ phát sinh, gia hạn):} ghi thêm đồ uống, in ấn, thiết bị vào đơn đang dùng; gia hạn thêm giờ nếu chỗ còn trống; thu tiền các khoản này bằng tiền mặt hoặc VietQR.
    \item \textbf{FR-19 (Bảo trì):} tạo lịch bảo trì cho một không gian; hệ thống xử lý các đơn bị ảnh hưởng (hủy và hoàn tiền hoặc hoàn phần thời gian bị mất).
    \item \textbf{FR-20 (Hủy đơn thay khách):} hủy đơn cho khách trong trường hợp ngoại lệ, bắt buộc nhập lý do, có thể miễn phí phạt; thao tác được ghi vào nhật ký kiểm toán.
\end{itemize}

\subsection{Quản lý chi nhánh}
\begin{itemize}
    \item \textbf{FR-21 (Tầng và không gian):} thêm, sửa, xóa tầng và không gian của chi nhánh; tải ảnh cho từng không gian.
    \item \textbf{FR-22 (Sơ đồ mặt bằng):} vẽ sơ đồ bằng công cụ kéo thả hoặc tải lên tệp SVG; gắn phần tử trên sơ đồ với không gian tương ứng.
    \item \textbf{FR-23 (Bảng giá chi nhánh):} đặt giá theo giờ, ngày, tuần, tháng cho từng loại không gian; giá chi nhánh được ưu tiên hơn giá chung.
    \item \textbf{FR-24 (Chính sách hủy chi nhánh):} cấu hình các bậc hoàn tiền riêng của chi nhánh.
    \item \textbf{FR-25 (Nhân viên):} xem, tạo tài khoản, khóa hoặc mở khóa nhân viên của chi nhánh.
    \item \textbf{FR-26 (Hoàn tiền):} xem các yêu cầu hoàn tiền của chi nhánh; sau khi trả tiền cho khách thì đánh dấu đã xử lý (tiền mặt, chuyển khoản hoặc voucher), hoặc từ chối kèm lý do.
    \item \textbf{FR-27 (Báo cáo chi nhánh):} xem doanh thu, số đơn, tỷ lệ hủy, cơ cấu doanh thu theo loại không gian và xuất CSV.
\end{itemize}

\subsection{Quản trị viên hệ thống}
\begin{itemize}
    \item \textbf{FR-28 (Chi nhánh):} tạo, sửa, tạm dừng chi nhánh; cấu hình địa chỉ và giờ mở cửa.
    \item \textbf{FR-29 (Danh mục dùng chung):} quản lý loại không gian, tiện ích và gắn tiện ích cho từng loại không gian.
    \item \textbf{FR-30 (Người dùng):} tra cứu tài khoản, đổi vai trò, gán chi nhánh, khóa tài khoản.
    \item \textbf{FR-31 (Giá và chính sách chung):} đặt bảng giá chung và chính sách hủy chung cho các chi nhánh chưa có cấu hình riêng.
    \item \textbf{FR-32 (Khuyến mãi và hạng thành viên):} tạo mã giảm giá; cấu hình bốn hạng thành viên (Bronze, Silver, Gold, Platinum) với ngưỡng chi tiêu, số đơn và mức giảm.
    \item \textbf{FR-33 (Báo cáo toàn hệ thống):} xem doanh thu toàn chuỗi và so sánh giữa các chi nhánh.
    \item \textbf{FR-34 (Nhật ký kiểm toán):} tìm kiếm nhật ký theo thời gian, người thực hiện, đối tượng và loại hành động.
\end{itemize}

\section{Yêu cầu phi chức năng}
\label{sec:yeucau_phichucnang}

\begin{enumerate}
    \item \textbf{NFR-01 (Hiệu năng):} các thao tác đọc (danh sách chi nhánh, sơ đồ, tình trạng chỗ) phản hồi dưới 500 ms; các thao tác ghi có kiểm tra tương tranh (tạo đơn, thanh toán, hủy) xử lý tại máy chủ dưới 1,5 giây, trong điều kiện mạng thông thường.
    \item \textbf{NFR-02 (Toàn vẹn dữ liệu):} không có hai đơn đang hiệu lực giao nhau trên cùng một chỗ, kể cả khi có yêu cầu đồng thời; với mọi đơn bị hủy, số tiền hoàn cộng phí phạt bằng tổng tiền đơn.
    \item \textbf{NFR-03 (Bảo mật):} giao tiếp qua HTTPS; mật khẩu băm bằng BCrypt; token nội bộ ký HS384 với khóa tối thiểu 384 bit; khóa bí mật được cấp qua biến môi trường, không ghi trong mã nguồn.
    \item \textbf{NFR-04 (Độ tin cậy):} các tác vụ định kỳ xử lý lỗi của từng đơn riêng lẻ, không làm dừng cả lượt xử lý; máy chủ tắt có kiểm soát (graceful shutdown) để không cắt ngang giao dịch đang chạy.
    \item \textbf{NFR-05 (Khả năng sử dụng):} giao diện tiếng Việt, dùng được trên màn hình rộng từ 375 px (điện thoại) đến 1920 px (máy tính), có chế độ sáng và tối.
    \item \textbf{NFR-06 (Khả năng bảo trì):} mã nguồn chia thành các tầng controller, service, repository; các luồng nghiệp vụ chính (đặt chỗ, thanh toán, hủy, check-in, phân quyền) có kiểm thử tự động.
\end{enumerate}

\section{Quy tắc nghiệp vụ}
\label{sec:quytac_nghiepvu}

Các quy tắc nghiệp vụ được đánh mã \mt{BR-xx} và được tham chiếu ở các chương sau (Bảng~\ref{tab:rule_concurrency} và Bảng~\ref{tab:rule_payment}).

\begin{table}[!htbp]
\centering
\caption{Quy tắc nghiệp vụ về đặt chỗ}
\label{tab:rule_concurrency}
\small
\renewcommand{\arraystretch}{1.2}
\begin{tabularx}{\linewidth}{|P{1.3cm}|L|P{4.6cm}|}
\hline
\textbf{Mã} & \textbf{Quy tắc} & \textbf{Mục đích} \\
\hline
BR-01 & Hai đơn ở trạng thái \mt{PENDING\_PAYMENT}, \mt{CONFIRMED} hoặc \mt{CHECKED\_IN} không được giao nhau về thời gian trên cùng một không gian; khoảng thời gian là nửa mở $[start, end)$ & Không đặt trùng chỗ \\
\hline
BR-02 & Đơn mới có hạn thanh toán 15 phút; quá hạn thì chuyển sang \mt{EXPIRED} và giải phóng chỗ & Tránh giữ chỗ ảo \\
\hline
BR-03 & Mỗi người dùng có tối đa 3 đơn \mt{PENDING\_PAYMENT} cùng lúc & Ngăn chiếm chỗ hàng loạt \\
\hline
BR-04 & Đơn theo giờ phải nằm trong giờ mở cửa của một ngày; đơn theo ngày, tuần, tháng phải bắt đầu trong giờ mở cửa; giờ tính theo múi giờ Việt Nam & Khách chỉ đến khi chi nhánh có người phục vụ \\
\hline
BR-05 & Không đặt cho thời điểm đã qua; khung giờ hiện tại vẫn được đặt & Dữ liệu lịch sử nhất quán \\
\hline
BR-06 & Số đơn vị tính tiền do máy chủ tính từ khoảng thời gian, làm tròn lên (giờ, ngày, tuần, tháng dương lịch); không dùng số đơn vị do trình duyệt gửi & Chống sửa giá từ phía trình duyệt \\
\hline
BR-07 & Thời lượng tối đa mỗi đơn: 24 giờ, 30 ngày, 52 tuần hoặc 12 tháng & Tránh nhập sai khoảng thời gian quá dài \\
\hline
BR-08 & Chi nhánh của đơn được suy ra từ không gian (không gian $\rightarrow$ tầng $\rightarrow$ chi nhánh), không nhận từ yêu cầu & Chống giả mạo chi nhánh \\
\hline
BR-09 & Tạo lịch bảo trì dùng chung khóa tư vấn với luồng đặt chỗ; không tạo lịch bảo trì trong quá khứ & Đặt chỗ và bảo trì không lọt qua nhau \\
\hline
BR-10 & Check-in được phép từ 30 phút trước giờ bắt đầu đến giờ kết thúc; gói nhiều ngày check-in lại được sau mỗi lần check-out & Kiểm soát giờ sử dụng \\
\hline
\end{tabularx}
\end{table}

\begin{table}[!htbp]
\centering
\caption{Quy tắc nghiệp vụ về thanh toán, hủy đơn và hoàn tiền}
\label{tab:rule_payment}
\small
\renewcommand{\arraystretch}{1.2}
\begin{tabularx}{\linewidth}{|P{1.3cm}|L|P{4.6cm}|}
\hline
\textbf{Mã} & \textbf{Quy tắc} & \textbf{Mục đích} \\
\hline
BR-11 & Khách chỉ hủy được đơn trước giờ bắt đầu; đơn \mt{CHECKED\_IN} không hủy được, khách muốn về sớm thì check-out & Không hủy sau khi đã dùng để đòi hoàn tiền \\
\hline
BR-12 & Chính sách hủy của chi nhánh được xét trước; nếu không bậc nào khớp thì xét chính sách chung; vẫn không khớp thì hoàn 0\%. Khoảng thời gian được tính theo phút và so với nửa khoảng $[min, max)$ & Chi nhánh tự chủ nhưng luôn có chính sách dự phòng; không chồng lấn tại biên \\
\hline
BR-13 & Số tiền hoàn cộng phí phạt bằng tổng tiền đơn; số tiền hoàn không vượt quá số tiền thực thu & Nhất quán tài chính \\
\hline
BR-14 & Hoàn tiền được đưa vào hàng chờ; quản lý chi nhánh hoặc quản trị viên trả tiền cho khách rồi đánh dấu đã xử lý, hoặc từ chối kèm lý do; hệ thống không tự chi tiền qua cổng thanh toán & Kiểm soát dòng tiền chi ra \\
\hline
BR-15 & Tiền về khi đơn đã \mt{EXPIRED} hoặc \mt{CANCELLED}, hoặc khách thanh toán trùng, thì tạo yêu cầu hoàn tiền thay vì xác nhận đơn & Bảo vệ quyền lợi khách khi thanh toán muộn \\
\hline
BR-16 & Khách đặt trên web chỉ thanh toán qua PayOS hoặc MoMo; tiền mặt chỉ do nhân viên xác nhận tại quầy & Tránh giữ chỗ trực tuyến rồi không đến trả tiền \\
\hline
BR-17 & Đơn \mt{CONFIRMED} hết giờ mà chưa từng check-in chuyển sang \mt{NO\_SHOW} và không được hoàn tiền & Chỗ đã được giữ trọn thời gian \\
\hline
BR-18 & Không check-out khi còn dịch vụ phát sinh chưa thanh toán; trả chỗ muộn quá 15 phút phải xử lý phí muộn trước & Không thất thu \\
\hline
BR-19 & Doanh thu chỉ tính trên đơn \mt{COMPLETED} và \mt{NO\_SHOW} theo thời điểm kết thúc, trừ dịch vụ chưa thu và các khoản đã hoàn & Không tính đơn chưa sử dụng \\
\hline
\end{tabularx}
\end{table}

\section{Sơ đồ use-case}
\label{sec:usecase_tongquat}

\subsection{Sơ đồ tổng quát}
Hình~\ref{fig:uc_tongquat} thể hiện ranh giới hệ thống và các nhóm chức năng của bốn tác nhân.

\begin{figure}[!htbp]
  \centering
  \includegraphics[width=\linewidth,height=0.72\textheight,keepaspectratio]{Chuong5/uc-tongquat.png}
  \caption{Sơ đồ use-case tổng quát của CoSpace}
  \label{fig:uc_tongquat}
\end{figure}

\subsection{Khách hàng}
Khách hàng tự thực hiện các thao tác: tìm chỗ trên sơ đồ, đặt chỗ, thanh toán, quản lý lịch sử đơn và kết nối đối tác (Hình~\ref{fig:uc_customer}).

\begin{figure}[!htbp]
  \centering
  \includegraphics[width=\linewidth,height=0.70\textheight,keepaspectratio]{Chuong5/uc-khachhang.png}
  \caption{Sơ đồ use-case của khách hàng}
  \label{fig:uc_customer}
\end{figure}

\subsection{Nhân viên quầy}
Nhân viên tạo đơn cho khách vãng lai, thu tiền mặt, check-in, check-out và ghi nhận dịch vụ phát sinh (Hình~\ref{fig:uc_staff}).

\begin{figure}[!htbp]
  \centering
  \includegraphics[width=\linewidth,height=0.70\textheight,keepaspectratio]{Chuong5/uc-nhanvien.png}
  \caption{Sơ đồ use-case của nhân viên quầy}
  \label{fig:uc_staff}
\end{figure}

\subsection{Quản lý chi nhánh}
Quản lý chi nhánh cấu hình mặt bằng, sơ đồ, bảng giá, chính sách hủy, quản lý nhân viên và xử lý hoàn tiền (Hình~\ref{fig:uc_branch}).

\begin{figure}[!htbp]
  \centering
  \includegraphics[width=\linewidth,height=0.70\textheight,keepaspectratio]{Chuong5/uc-quanlychinhanh.png}
  \caption{Sơ đồ use-case của quản lý chi nhánh}
  \label{fig:uc_branch}
\end{figure}

\subsection{Quản trị viên hệ thống}
Quản trị viên quản lý chi nhánh, người dùng, chính sách chung, xem báo cáo toàn hệ thống và nhật ký kiểm toán (Hình~\ref{fig:uc_admin}).

\begin{figure}[!htbp]
  \centering
  \includegraphics[width=\linewidth,height=0.70\textheight,keepaspectratio]{Chuong5/uc-quantrihethong.png}
  \caption{Sơ đồ use-case của quản trị viên hệ thống}
  \label{fig:uc_admin}
\end{figure}

\section{Ma trận phân quyền}
\label{sec:ma_tran_phan_quyen}

Bảng~\ref{tab:matran_phanquyen} tóm tắt quyền của từng vai trò trên các nhóm tài nguyên. Ký hiệu: C (tạo), R (xem), U (cập nhật), D (xóa hoặc ngừng sử dụng); ``CN'' nghĩa là chỉ trong chi nhánh của người dùng.

\begin{table}[!htbp]
\centering
\caption{Ma trận phân quyền theo vai trò}
\label{tab:matran_phanquyen}
\small
\renewcommand{\arraystretch}{1.15}
\begin{tabularx}{\linewidth}{|L|P{2.0cm}|P{2.2cm}|P{2.2cm}|P{2.0cm}|}
\hline
\textbf{Tài nguyên} & \textbf{Khách hàng} & \textbf{Nhân viên} & \textbf{Quản lý CN} & \textbf{Quản trị} \\
\hline
Chi nhánh, không gian (xem) & R & R (CN) & R (CN) & CRUD \\
\hline
Tầng, không gian, sơ đồ & -- & -- & CRUD (CN) & CRUD \\
\hline
Đơn đặt chỗ của mình & CR, hủy & -- & -- & -- \\
\hline
Đơn tại quầy & -- & CR (CN) & CR (CN) & CR \\
\hline
Xác nhận tiền mặt & -- & U (CN) & U (CN) & U \\
\hline
Check-in, check-out & -- & U (CN) & U (CN) & U \\
\hline
Dịch vụ phát sinh trên đơn & C (đơn của mình) & CRU (CN) & CRU (CN) & CRU \\
\hline
Lịch bảo trì & -- & CRU (CN) & CRU (CN) & CRU \\
\hline
Bảng giá, chính sách hủy & R & -- & CRUD (CN) & CRUD (chung) \\
\hline
Hàng chờ hoàn tiền & -- & -- & RU (CN) & RU \\
\hline
Tài khoản nhân viên & -- & -- & CRU (CN) & CRUD \\
\hline
Báo cáo, xuất CSV & -- & R (thống kê trong ngày) & R (CN) & R \\
\hline
Hồ sơ, kết nối, bảng tin & CRUD (của mình) & -- & -- & -- \\
\hline
Trợ lý ảo & Sử dụng & -- & -- & -- \\
\hline
Khuyến mãi, hạng thành viên & R & -- & -- & CRUD \\
\hline
Nhật ký kiểm toán & -- & -- & -- & R \\
\hline
\end{tabularx}
\end{table}
